% thm-sharp-radii-lil -- CERW.Frozen.sharp_radii_lil (draft). Source: limit-shapes.tex:165-170, label thm:sharp (part (ii), eq:planar-lil-lower)
\paragraph{Paper (verbatim).} Theorem environment (lines 156--188); this node is part (ii), lines 165--170:
\begin{verbatim}
\begin{theorem}[Sharpness]\label{thm:sharp}
Let $d\geq2$ and $0<\drift<\nf{1}{d}$, and consider the centrally excited random walk.
\begin{enumerate}[label=\textup{(\roman*)}]
\item \underline{\emph{Radii in the plane}}: let $d=2$. For every $p>0$ there exists $c(\drift,p)>0$ such that, for all sufficiently large~$n$,
\begin{equation}\label{eq:planar-polynomial-lower}
\P\bigl(r_n-\Rin(n)\geq c\sqrt{r_n\log n}\bigr)\geq n^{-p}
\qquad\text{and}\qquad
\P\bigl(\Rout(n)-r_n\geq c\sqrt{r_n\log n}\bigr)\geq n^{-p}\, .
\end{equation}
\item \underline{\emph{Iterated logarithm for the radii in the plane}}: let $d=2$. Almost surely,
\begin{equation}\label{eq:planar-lil-lower}
\limsup_{n\to\infty}\frac{r_n-\Rin(n)}{\sqrt{r_n\log\log n}}\geq\frac1{\sqrt{10\pi\drift}}
\qquad\text{and}\qquad
\limsup_{n\to\infty}\frac{\Rout(n)-r_n}{\sqrt{r_n\log\log n}}\geq\frac1{\sqrt{10\pi\drift}}\, .
\end{equation}
\item \underline{\emph{Local times near the origin}}: there exists $c(d,\drift)>0$ such that, for all sufficiently large~$n$, with probability at least $1-n^{-c}$,
\begin{equation}\label{eq:sharp-bulk}
\max_{\substack{x\in\Z^d\\|x|\leq r_n^{\nf12}}}\bigl|\ell_n(x)-2d\drift(r_n-|x|)\bigr|
\geq c\begin{cases}\sqrt{r_n}\log n,&d=2,\\\sqrt{r_n\log n},&d\geq3\end{cases}\, .
\end{equation}
\item \underline{\emph{Difference of the radii in the plane}}: let $d=2$.
\begin{enumerate}[label=\textup{(\alph*)}]
\item For every $p>0$ there exists $c(\drift,p)>0$ such that, for all sufficiently large~$n$,
\begin{equation}\label{eq:planar-width-polynomial}
\P\bigl(\Rout(n)-\Rin(n)\geq c\sqrt{r_n\log n}\bigr)\geq n^{-p}\, .
\end{equation}
\item Almost surely,
\begin{equation}\label{eq:planar-width-lil}
\limsup_{n\to\infty}\frac{\Rout(n)-\Rin(n)}{\sqrt{r_n\log\log n}}\geq\sqrt{\frac{\pi}{3\drift}}\, .
\end{equation}
\end{enumerate}
\end{enumerate}
\end{theorem}
\end{verbatim}
\paragraph{Transcription.} $d=2$ is \texttt{hd : d = 2}; the external Stout law of the iterated logarithm is carried as the hypothesis
\texttt{hLIL : CERW.External.StoutLIL.{u}} (author ruling: the proof reaches Stout's theorem through
Theorem~\ref{thm:moment-fluctuations}(ii), line 1410). ``Almost surely $\limsup_n a_n\ge L$'' is written junk-free
(repository convention) as \texttt{∀ᵐ ω, ∀ δ > 0, ∃ᶠ n, (L - δ) * √(r n * log (log n)) ≤ a n} with
$L=1/\sqrt{10\pi\kappa}$, $a_n=r_n-R_{\rm in}(n)$ for the first display and $a_n=R_{\rm out}(n)-r_n$ for the second; both
displays are conjuncts under the same \texttt{∀ δ} inside one \texttt{∀ᵐ}. Dividing by $\sqrt{r_n\log\log n}>0$ (true for $n\ge3$)
recovers the ratio form, and the finitely many $n$ with $\log\log n\le0$ do not affect a ``frequently'' statement.
